\chapter{Métodos numéricos iterativos}
\label{cap:metodos}

Este capítulo presenta los \datoNMetodos{} métodos que se comparan en el trabajo.
Para cada uno se enuncia la iteración, se establece su orden de convergencia
teórico y se analiza \emph{qué le ocurre sobre el problema concreto} del \cref{cap:estructura}. Este último punto es el que da contenido al capítulo: sobre
una función lineal a trozos, varios métodos pueden terminar en pocos pasos, y
su comportamiento se explica mejor por la estructura del problema que por el
orden asintótico que les asigna la teoría genérica.

Los métodos se agrupan en tres familias según la información que utilizan y la
garantía que ofrecen.

\begin{table}[H]
\centering
\caption[Taxonomía de los métodos]{Taxonomía de los métodos comparados. La columna
$c$ indica el costo nominal por iteración; para los métodos analíticos supone
una obtención conjunta de valor y derivada. $E=p^{1/c}$ es el índice de
eficiencia de Ostrowski (\cref{def:ostrowski}); el costo efectivo incluye todas
las llamadas al motor.}
\label{tab:taxonomia}
\small
\begin{tabular}{@{}llccc l@{}}
\toprule
\textbf{Familia} & \textbf{Método} & \textbf{$p$} & \textbf{$c$} & \textbf{$E$} & \textbf{Garantía} \\
\midrule
\multirow{8}{*}{Encaje}
 & Bisección          & $1$      & $1$ & $1{,}000$ & global, solo continuidad \\
 & Regula falsi       & $1$      & $1$ & $1{,}000$ & global \\
 & Illinois           & $1{,}442$& $1$ & $1{,}442$ & global \\
 & Pegasus            & $1{,}642$& $1$ & $1{,}642$ & global \\
 & Anderson--Björck   & $1{,}710$& $1$ & $1{,}710$ & global \\
 & Ridders            & $2$      & $2$ & $1{,}414$ & global \\
 & Brent              & superlineal& $1$ & --- & global, con encaje \\
 & ITP                & $1{,}618$& $1$ & $1{,}618$ & global, peor caso óptimo \\
\midrule
\multirow{9}{*}{Abiertos}
 & Secante            & $1{,}618$& $1$ & $1{,}618$ & local \\
 & Newton (dif.\ fin.)& $2$      & $3$ & $1{,}260$ & local \\
 & Newton (analítico) & $2$      & $1$ & $2{,}000$ & local \\
 & Steffensen         & $2$      & $2$ & $1{,}414$ & local \\
 & Steffensen escalado& $2$      & $2$ & $1{,}414$ & local \\
 & Müller             & $1{,}839$& $1$ & $1{,}839$ & local \\
 & Halley             & $3$      & $3$ & $1{,}442$ & local \\
 & Punto fijo relajado& $1$      & $1$ & $1{,}000$ & local, contracción \\
 & Aitken--Steffensen & $2$      & $2$ & $1{,}414$ & local \\
\midrule
\multirow{2}{*}{No suaves}
 & Newton semisuave   & $2$      & $1$ & $2{,}000$ & local, $BD$-regularidad \\
 & Newton salvaguardado& finito local& $1$ & --- & global; bisección periódica \\
\bottomrule
\end{tabular}
\end{table}

\section{Marco general y criterios de comparación}

\subsection{El costo relevante}

\begin{remark}[Por qué se cuentan evaluaciones y no iteraciones]
\label{rem:costo}
Una llamada al motor construye la proyección actuarial de $\datoNMeses$
periodos y domina el costo de las pocas operaciones escalares de la iteración.
La medida principal de trabajo es, por ello, el número de llamadas efectivas
al motor, junto con el tiempo medido. Deben incluirse las llamadas que
obtienen derivadas y las que preparan parámetros de relajación. El costo
nominal de Newton analítico es una llamada por punto solo cuando valor y
derivada se calculan conjuntamente y se reutilizan; si se calculan por
separado, ambas llamadas se contabilizan. La secante necesita un nuevo valor
por paso y Newton con diferencias centradas requiere, en general, tres.
\end{remark}

La instrumentación memoiza los valores ya calculados por punto. Una consulta
resuelta desde esa memoria no constituye una nueva ejecución del motor. La
misma convención debe aplicarse a todos los métodos y a la preparación que
cada uno necesita.

\subsection{Criterios de parada y el desajuste dimensional}
\label{sec:parada}

La tolerancia del residuo se expresa en unidades de capital. En los métodos
con encaje también puede certificarse el error mediante el ancho del intervalo:
\begin{equation}
\label{eq:parada}
  |f(x_k)|<\varepsilon_f
  \qquad\text{o, con encaje válido,}\qquad
  \frac{b_k-a_k}{2}\le\varepsilon_x ,
\end{equation}
devolviendo en el segundo caso el punto medio del intervalo. Un paso pequeño
$|x_k-x_{k-1}|$ se utiliza como diagnóstico de estancamiento; por sí solo no
certifica una raíz, pues un recorte al dominio puede producir un paso nulo
con residuo no nulo.

\begin{criticism}[Conversión del residuo en error de mesada]
\label{crit:parada}
El motor de producción utiliza una tolerancia de
$\varepsilon_f=\datoEpsilonProd$ peso de capital. Para relacionarla con la
exactitud de la mesada, supóngase que las pendientes de $\VP$ en el intervalo
que une $x$ con $M^*$ están acotadas inferiormente por $m>0$. Entonces
\[
  |x-M^*|\le\frac{|f(x)|}{m}.
\]
En un mismo tramo afín puede tomarse $m=\VP'$, de modo que la pendiente
$\datoPendienteAfin$ convierte un residuo de un peso en un error de mesada del
orden de $4\cdot10^{-3}$ pesos. Para garantizar una tolerancia de mesada
$\varepsilon$, es suficiente exigir $|f(x)|\le m\varepsilon$. Si se atraviesan
quiebres, la derivada en un único punto no sustituye necesariamente a esa cota
inferior. La recomendación~R3 de la \cref{sec:hallazgos-codigo} consiste en
declarar las unidades de la tolerancia y certificar su conversión mediante
una cota válida o mediante el encaje.
\end{criticism}

\section{Métodos con encaje}

Todos comparten la hipótesis inicial $f(a)f(b)<0$ y la propiedad invariante de
mantener un intervalo que contiene una raíz. Su virtud común es que la garantía de
convergencia \emph{no usa diferenciabilidad}, y por tanto sobrevive intacta a la
estructura no suave de $\VP$.

\subsection{Bisección}

\begin{algorithm}[H]
\caption{Bisección}
\label{alg:biseccion}
\SetAlgoLined
\KwIn{$f$, $[a,b]$ con $f(a)f(b)<0$, tolerancias}
\While{no se cumpla \cref{eq:parada}}{
  $x\leftarrow (a+b)/2$\;
  \eIf{$f(a)f(x)<0$}{$b\leftarrow x$}{$a\leftarrow x$, $f(a)\leftarrow f(x)$}
}
\KwOut{$x$}
\end{algorithm}

Su análisis es el \cref{thm:cota-biseccion}. Sobre el problema de este trabajo, con
$[a,b]=[S_0,25S_0]$ y $\varepsilon_x=10^{-8}$, la cota da
$\lceil\log_2(24S_0/10^{-8})\rceil\approx52$ iteraciones. Con la tolerancia de
producción en el residuo el número observado es $\datoEvBis$ evaluaciones en
promedio\footnote{La discrepancia entre la cota teórica de unas cincuenta y dos iteraciones y las treinta y tres observadas se explica porque el criterio de parada efectivo es el del residuo y no el del paso: con una pendiente de $\datoPendienteAfin$, un residuo menor que un peso corresponde a un error en la mesada del orden de $4\cdot10^{-3}$ pesos, muy por encima de la tolerancia $\varepsilon_x=10^{-8}$ que da lugar a la cota.}.

\begin{remark}[Una advertencia metodológica]
\label{rem:sesgo-diadico}
Un detalle del diseño experimental resulta crucial. Si la raíz se sitúa en un
múltiplo «redondo» del salario mínimo, los puntos medios sucesivos de la bisección
pueden coincidir \emph{exactamente} con ella por pura aritmética del encaje. En una
versión preliminar de los experimentos, con objetivos en $2{,}5$ y $5{,}0$
salarios mínimos, la bisección aparentaba converger en $6$ evaluaciones ---
compitiendo con Newton --- por este artefacto. Los escenarios definitivos usan
multiplicadores irracionales (por ejemplo, $1+\sqrt2$ y $3+\sqrt5$) para
reducir ese sesgo. El efecto sobre la conclusión es grande: el costo medio
de la bisección pasa de $21{,}6$ a $\datoEvBis$ evaluaciones\footnote{Se documenta este episodio porque ilustra un riesgo metodológico general en la comparación de algoritmos: los conjuntos de prueba construidos con valores «redondos» pueden favorecer sistemáticamente a los métodos cuya aritmética coincide con esa forma de redondeo. La bisección, que divide en potencias de dos, es especialmente vulnerable a ello.}.
\end{remark}

\subsection{Regula falsi y sus correcciones}

La regula falsi sustituye el punto medio por la raíz de la cuerda:
\begin{equation}
\label{eq:regula}
  x \;=\; b - f(b)\,\frac{b-a}{f(b)-f(a)} .
\end{equation}

\begin{proposition}[Estancamiento de la regula falsi]
\label{prop:estancamiento}
Si $f$ es continua, estrictamente convexa y estrictamente creciente en $[a,b]$,
con $f(a)<0<f(b)$, el extremo $b$ se conserva en las iteraciones exactas
de regula falsi y el ancho del encaje tiende a $b-x^*>0$. Si además $f$ es
continuamente diferenciable en un entorno de la raíz y $f'(x^*)>0$, la
convergencia es lineal con razón
$R=1-f'(x^*)(b-x^*)/f(b)\in(0,1)$.
\end{proposition}
\begin{proof}
La convexidad estricta sitúa la gráfica estrictamente por debajo de la
cuerda en el interior. En la raíz de la cuerda se tiene, por tanto,
$f(x)<0$ y se actualiza $a\leftarrow x$. El extremo derecho se conserva.
La sucesión de extremos izquierdos crece hacia la raíz y su límite da
el ancho indicado. La linealización de \cref{eq:regula} respecto del
extremo izquierdo en $x^*$ proporciona $R$. La pendiente de la cuerda
entre $x^*$ y $b$ es estrictamente mayor que $f'(x^*)$, de donde
$0<R<1$. Con convexidad no estricta puede haber tramos afines y
terminación exacta; ese caso no implica estancamiento asintótico.
\end{proof}

En el modelo afín a trozos, la conservación de un extremo puede coexistir
con una buena reducción del residuo. Deben distinguirse el ancho del
encaje y el error del punto propuesto; la proposición anterior no se aplica
como una hipótesis de convexidad estricta global a $\VP$.

Las tres correcciones clásicas atacan el estancamiento reduciendo por un factor
$\gamma$ el valor de $f$ en el extremo retenido:
\begin{align*}
  \text{Illinois:}&\quad \gamma=\tfrac12,\\
  \text{Pegasus:}&\quad \gamma=\frac{f(b)}{f(b)+f(x)},\\
  \text{Anderson--Björck:}&\quad \gamma=\max\Bigl\{1-\frac{f(x)}{f(b)},\ \tfrac12\Bigr\}.
\end{align*}
Los órdenes resultantes son $1{,}442$, $1{,}642$ y $1{,}710$ respectivamente
\citep{dowell1972,anderson1973}, con una sola evaluación por paso: son, en
términos del índice de Ostrowski, de los métodos más eficientes del catálogo entre
los que conservan garantía global.

\begin{remark}
La coincidencia de conteos medios entre variantes no implica que sus
sucesiones de iterados sean idénticas. La estructura afín explica que las
diferencias puedan ser pequeñas: cuando los dos puntos que definen la
cuerda pertenecen al mismo tramo que contiene la raíz, la interpolación
lineal es exacta. Fuera de esa situación deben contrastarse las
trayectorias, además de sus conteos agregados.
\end{remark}

\subsection{Ridders}

Supone que $g(x)=f(x)e^{Q(x-x_1)}$ es lineal para un $Q$ adecuado. Evaluando en el
punto medio $x_2$ y despejando se obtiene
\begin{equation}
\label{eq:ridders}
  x_3 \;=\; x_2 + (x_2-a)\,\frac{\sgn\bigl(f(a)-f(b)\bigr)\,f(x_2)}
                                {\sqrt{f(x_2)^2-f(a)f(b)}} .
\end{equation}
El iterado $x_3$ siempre cae dentro de $[a,b]$, de modo que conserva el encaje, y
la convergencia es cuadrática por \emph{par} de evaluaciones, es decir de orden
$\sqrt2$ por evaluación \citep{ridders1979}.

\subsection{Brent}

Combina interpolación cuadrática inversa, secante y bisección, aceptando el
candidato interpolado solo si satisface condiciones de progreso
\citep{brent1973}. Conserva el encaje y presenta convergencia superlineal en
el caso regular. La interpolación no implica que el número de iteraciones
coincida con la cota de la bisección en todo problema. Se utiliza como
referencia de un método híbrido disponible en bibliotecas científicas, como
\texttt{scipy.optimize.brentq}.

\subsection{ITP}

El método \emph{Interpolate--Truncate--Project} \citep{oliveira2021} es el
resultado más reciente del catálogo. En cada iteración:
\begin{enumerate}[nosep,leftmargin=*]
  \item \textbf{Interpolar:} calcular el candidato de regula falsi $x_f$.
  \item \textbf{Truncar:} desplazarlo hacia el punto medio a lo sumo
        $\delta=\kappa_1|b-a|^{\kappa_2}$, obteniendo $x_t$.
  \item \textbf{Proyectar:} proyectar $x_t$ sobre el intervalo minmax
        $[x_{1/2}-r,\;x_{1/2}+r]$ con
        $r=\varepsilon 2^{n_{\max}-j}-\tfrac12(b-a)$.
\end{enumerate}
La etapa 3 es la que garantiza que ITP nunca gaste más de $n_0$ iteraciones
adicionales sobre la cota de la bisección\footnote{El resultado de \citet{oliveira2021} es notable porque durante décadas se consideró que la optimalidad en el peor caso y la convergencia superlineal eran objetivos incompatibles: los métodos rápidos degradaban en el peor caso y los robustos no aceleraban en el caso bueno. ITP alcanza ambas.}, y --- punto clave para este trabajo ---
esa garantía no usa diferenciabilidad, luego sobrevive en los quiebres.

\begin{remark}[Por qué ITP resulta caro aquí]
Los experimentos muestran que ITP requiere $\datoEvITP$ evaluaciones, muy por
encima de Brent. No es un defecto del método sino del ajuste de sus parámetros: la
etapa de proyección está calibrada para \emph{garantizar} el peor caso, y con el
intervalo de partida $[S_0,25S_0]$, muy amplio respecto de la raíz típica, la
proyección domina sobre la interpolación durante muchas iteraciones. Cabe
esperar que un intervalo inicial más ajustado reduzca ese costo, aunque esa
variante no se midió en este trabajo. Es un ejemplo instructivo de que la optimalidad en el peor caso se paga en el caso
medio.
\end{remark}

\section{Métodos abiertos}

\subsection{Secante}

\begin{equation}
\label{eq:secante}
  x_{k+1} \;=\; x_k - f(x_k)\,\frac{x_k-x_{k-1}}{f(x_k)-f(x_{k-1})} .
\end{equation}

La secante es el método configurado por defecto en el motor de valoración.

\begin{theorem}[Orden de la secante]
\label{thm:orden-secante}
Sea $f\in C^2$ en un entorno de una raíz simple $x^*$, con
$f''(x^*)\ne0$. Para una sucesión de secante bien definida, iniciada con
dos puntos distintos suficientemente próximos a $x^*$, que converge a
$x^*$ sin terminación finita, los errores absolutos satisfacen
\[
 e_{k+1}=C\,e_ke_{k-1}(1+o(1)),\qquad
 C=\frac{|f''(x^*)|}{2|f'(x^*)|}>0.
\]
Su orden es $\varphi=(1+\sqrt5)/2\approx1{,}618$. La condición de no
degeneración excluye, en particular, el tramo afín, en el que puede haber
terminación exacta en lugar de un régimen asintótico.
\end{theorem}

\begin{proposition}[La secante sobre una función afín a trozos]
\label{prop:secante-afin}
Si $x_{k-1}$ y $x_k$ son distintos, pertenecen al mismo tramo afín de $f$
con pendiente no nula, y ese tramo contiene la raíz, entonces
$x_{k+1}=x^*$ exactamente en aritmética exacta.
\end{proposition}
\begin{proof}
En ese tramo $f(u)=\alpha u+\beta$ con $\alpha\neq0$. La pendiente de la secante es
$(f(x_k)-f(x_{k-1}))/(x_k-x_{k-1})=\alpha$, y \cref{eq:secante} da
$x_{k+1}=x_k-f(x_k)/\alpha=-\beta/\alpha=x^*$.
\end{proof}

Esta propiedad explica el buen comportamiento observado de la secante. Cuando
la sucesión alcanza la solución en pocos pasos, una regresión sobre los errores
es un descriptor de esa sucesión y no una estimación fiable del orden
asintótico.

\subsection{Newton con diferencias finitas}

Sustituye la derivada por una diferencia centrada:
\begin{equation}
\label{eq:newton-fd}
  x_{k+1} \;=\; x_k - \frac{f(x_k)}{D_h f(x_k)},
  \qquad
  D_hf(x)=\frac{f(x+h)-f(x-h)}{2h} .
\end{equation}

\begin{proposition}[Exactitud y efecto de los quiebres en la diferencia centrada]
\label{prop:sesgo-diferencias}
Sea $f$ continua y lineal a trozos, con conjunto finito de quiebres
$\mathcal B$, y sea $h>0$.
\begin{enumerate}[label=(\alph*),leftmargin=*]
  \item Si $[x-h,x+h]\cap\mathcal B=\emptyset$, entonces
        $D_hf(x)=f'(x)$ en aritmética exacta.
  \item Si $\beta\in(x-h,x+h)$ es el único quiebre de la ventana y
        $s^-,s^+$ son las pendientes a sus lados, entonces
        \[
          D_hf(x)=\theta s^-+(1-\theta)s^+,
          \qquad \theta=\frac{\beta-x+h}{2h}\in(0,1).
        \]
        Si $x=\beta$, esta cantidad pertenece a $\dC f(\beta)$.
        Si $x\ne\beta$, no tiene por qué pertenecer a $\dC f(x)$.
  \item Si $\widetilde f(x\pm h)=f(x\pm h)+e_\pm$, con
        $|e_\pm|\le\eta_\pm$, el error propagado a la diferencia,
        considerando exactas la resta y la división finales, satisface
        \[
          \left|\frac{\widetilde f(x+h)-\widetilde f(x-h)}{2h}
                 -D_hf(x)\right|
          \le\frac{\eta_++\eta_-}{2h}.
        \]
\end{enumerate}
\end{proposition}
\begin{proof}
En (a), $f$ es afín en toda la ventana. Para (b),
\[
 f(x+h)-f(x-h)=s^-(\beta-x+h)+s^+(x+h-\beta),
\]
y basta dividir entre $2h$. Si $x=\beta$, la diferencia es
$(s^-+s^+)/2$ y pertenece al intervalo de Clarke. Si $x$ no es un quiebre,
$\dC f(x)=\{f'(x)\}$, que en general no contiene esa pendiente media.
Para (c), se resta la diferencia exacta de la calculada y se aplica
la desigualdad triangular.
\end{proof}

\begin{aporte}
\textbf{Elección del incremento.} El incremento $h$ equilibra la amplificación
del error de evaluación al dividir entre $2h$ y la mezcla de pendientes cuando
la ventana atraviesa un quiebre. La distancia desde $x$ al quiebre más próximo
determina si se produce el segundo efecto.

Aquí $f=\VP-\KA$. Cerca de la raíz puede haber cancelación entre dos cantidades
grandes; no se supone un error relativo pequeño respecto del residuo $f(x)$.
La estimación de $\eta_\pm$ debe incorporar las escalas de $\VP$ y de $\KA$,
además del error aritmético de las operaciones finales. El protocolo utiliza
$h=1$ peso. Las identidades anteriores no establecen por sí solas una ventana
universal de incrementos admisibles; su elección se contrasta en las regiones
y escalas efectivamente estudiadas.
\end{aporte}

La desventaja de Newton con diferencias finitas está en el costo y no en la exactitud: tres
llamadas al motor por iteración, frente a una de la secante. Con los órdenes
clásicos, el índice de Ostrowski es $2^{1/3}\approx1{,}260$ contra $1{,}618$; como
se advierte después de la \cref{def:ostrowski}, en este problema esos valores no
se interpretan como eficiencias asintóticas observadas. Lo que el experimento sí
refleja es el costo por iteración: $\datoEvNFD$ llamadas medias frente a
$\datoEvSec$ de la secante.

\subsection{Newton con derivada analítica y Newton semisuave}

Disponiendo de la expresión cerrada del \cref{thm:subdiferencial}, la iteración es
\begin{equation}
\label{eq:newton-semisuave}
  x_{k+1} \;=\; x_k - \frac{f(x_k)}{V_k},
  \qquad V_k\in\dB f(x_k),
\end{equation}
tomando $V_k=\VP'_+(x_k)$. Fuera de $\mathcal{B}$ coincide con el método de Newton
clásico; en un quiebre, es el método de Newton semisuave de Qi--Sun.

\begin{aporte}
\begin{theorem}[Terminación en un tramo afín y cota bajo convexidad]
\label{thm:terminacion-finita}
Sea $f$ continua, lineal a trozos y estrictamente creciente en
$[\alpha,\beta]$, con raíz $x^*$. En cada iteración, sea $V_k>0$ la pendiente
de uno de los tramos afines cuya clausura contiene $x_k$, y considérese
\cref{eq:newton-semisuave}. Entonces:
\begin{enumerate}[label=(\roman*),leftmargin=*]
  \item Si la clausura del tramo cuya pendiente se selecciona contiene
        también a $x^*$, entonces $x_{k+1}=x^*$ en aritmética exacta.
  \item Si además $f$ es convexa y $x_0>x^*$, los iterados anteriores a la
        terminación satisfacen $x^*\le x_{k+1}<x_k$. En cada paso que no alcanza
        la raíz se abandona el tramo seleccionado hacia otro situado
        estrictamente a su izquierda. El método termina en un número de
        pasos no mayor que el número de tramos cuyas clausuras intersectan
        $[x^*,x_0]$.
\end{enumerate}
La cota del apartado (ii) es una cota superior, no una igualdad.
\end{theorem}
\begin{proof}
En el tramo seleccionado $f(u)=a_ju+b_j$, con $a_j>0$. Por continuidad,
esta expresión vale en su clausura. Si esta contiene $x^*$, entonces
$x^*=-b_j/a_j$ y el paso proporciona exactamente ese valor.

Para (ii), mientras $x_k>x^*$ se tiene $f(x_k)>0$ y
$x_{k+1}=x_k-f(x_k)/V_k<x_k$. Por convexidad, la prolongación afín es una recta
de soporte: $f(u)\ge f(x_k)+V_k(u-x_k)$. Al evaluar en $u=x^*$ se obtiene
$x^*\le x_{k+1}$. Si $x_{k+1}$ permaneciera en la clausura del mismo tramo,
su valor de $f$ sería cero y el método habría terminado. En caso contrario,
está estrictamente a su izquierda. La finitud de los tramos da la cota.
\end{proof}
\end{aporte}

\begin{corollary}[Consecuencia para el problema de cotización]
\label{cor:un-paso}
Si el punto inicial $x_0$ y la solución $M^*$ pertenecen ambos al intervalo
abierto $I=(\rho_{n_{\max}}S_0,5S_0)$ del \cref{cor:region-afin}, el método
\cref{eq:newton-semisuave}, con la pendiente exacta de ese tramo, alcanza
$M^*$ en una sola iteración en aritmética exacta. Se requieren dos
evaluaciones de $f$ para calcular y comprobar el paso. El costo de obtener
la derivada debe incluirse en la contabilidad del motor.
\end{corollary}

El \cref{tab:terminacion} verifica el \cref{cor:un-paso} directamente.

\begin{table}[H]
\centering
\caption[Verificación de la terminación finita]{Verificación del
\cref{thm:terminacion-finita} sobre el escenario \texttt{A-R3}, cuya raíz es
$M^*=\datoMesadaEjemplo$ ($2{,}414$ salarios mínimos, dentro del tramo afín
$[1{,}15;\,5]\,S_0$). Los arranques interiores verifican la terminación en un
paso. Los demás describen el comportamiento observado y no establecen una
igualdad general entre el número de iteraciones y el de tramos atravesados.}
\label{tab:terminacion}
\small
\begin{tabular}{@{}lccc l@{}}
\toprule
$x_0/S_0$ & Tramo de $x_0$ & Iteraciones & Evaluaciones & Error final (pesos) \\
\midrule
$1{,}30$ & mismo que $M^*$ & $1$ & $2$ & $0$ \\
$2{,}00$ & mismo que $M^*$ & $1$ & $2$ & $0$ \\
$3{,}00$ & mismo que $M^*$ & $1$ & $2$ & $9{,}3\cdot10^{-10}$ \\
$4{,}50$ & mismo que $M^*$ & $1$ & $2$ & $1{,}9\cdot10^{-9}$ \\
$4{,}90$ & mismo que $M^*$ & $1$ & $2$ & $0$ \\
\midrule
$6{,}00$ & cruza racimo II & $2$ & $3$ & $9{,}3\cdot10^{-10}$ \\
$8{,}00$ & cruza racimo II & $2$ & $3$ & $0$ \\
$13{,}00$& cruza racimos II--III & $2$ & $3$ & $0$ \\
$20{,}00$& cruza racimos II--III & $2$ & $3$ & $0$ \\
$24{,}00$& cruza racimos II--III & $2$ & $3$ & $9{,}3\cdot10^{-10}$ \\
\bottomrule
\end{tabular}
\end{table}

El error residual de $9{,}3\cdot10^{-10}$ pesos es del orden de unas pocas unidades
en el último lugar (\emph{ulp}) de la representación de punto flotante de doble precisión
de un número del orden de $4\cdot10^6$: la solución alcanza el límite de la
aritmética disponible.

\subsection{Newton salvaguardado}

\begin{algorithm}[H]
\caption{Newton semisuave con salvaguarda de encaje}
\label{alg:newton-salvaguardado}
\SetAlgoLined
\KwIn{$f$, subgradiente $V(\cdot)$, $[a,b]$ con $f(a)f(b)<0$}
evaluar los extremos y devolver uno si satisface el criterio del residuo\;
$x\leftarrow(a+b)/2$, $c\leftarrow0$\;
\While{no se haya certificado la solución}{
  evaluar $f(x)$\;
  \If{$f(x)=0$ \KwSty{o} $|f(x)|<\varepsilon_f$}{devolver $x$\;}
  actualizar $[a,b]$ con $x$, conservando el cambio de signo\;
  \If{$(b-a)/2\le\varepsilon_x$}{
     evaluar y devolver $(a+b)/2$\;
  }
  \eIf{$c=2$}{
     $x\leftarrow(a+b)/2$, $c\leftarrow0$ \tcp*{bisección periódica}
  }{
     $v\leftarrow V(x)$\;
     \eIf{$v=0$ \KwSty{o} $x-f(x)/v\notin(a,b)$}{
        $x\leftarrow(a+b)/2$, $c\leftarrow0$\;
     }{
        $x\leftarrow x-f(x)/v$, $c\leftarrow c+1$\;
     }
  }
}
\end{algorithm}

\begin{theorem}[Convergencia global y terminación finita]
\label{thm:salvaguardado}
Sea $f$ continua y estrictamente creciente en el intervalo inicial
$[a_0,b_0]$, con $f(a_0)<0<f(b_0)$. El
\cref{alg:newton-salvaguardado} conserva el encaje y converge a su raíz única
$M^*$. Cada bloque de tres propuestas que no termina contiene al menos una
bisección; por tanto, después de evaluar e incorporar al encaje $k$
propuestas posteriores a la evaluación inicial, el ancho está acotado por
\[
  b_k-a_k\le (b_0-a_0)\,2^{-\lfloor k/3\rfloor}.
\]
La terminación por ancho devuelve un punto con error no mayor que
$(b_k-a_k)/2$. Esta cota de trabajo no es la cota por iteración de la
bisección pura.

Si además $f$ es lineal a trozos con finitos quiebres y el subgradiente
seleccionado es una pendiente lateral positiva en un entorno de $M^*$,
la iteración ideal, continuada hasta $f(x)=0$, alcanza la raíz en un
número finito de pasos en aritmética exacta. Con tolerancias positivas
puede detenerse antes mediante los certificados de \cref{eq:parada}.
En particular, un paso de Newton aceptado cuya pendiente corresponde a un
tramo cuya clausura contiene $M^*$ alcanza la raíz exactamente.
\end{theorem}
\begin{proof}
Cada evaluación que no termina reemplaza un extremo por un punto interior,
sin perder el cambio de signo. El contador permite a lo sumo dos propuestas
consecutivas de Newton antes de imponer un punto medio; un candidato inválido
también produce bisección. Cada una de estas reduce el ancho a la mitad y
los demás pasos no lo aumentan, lo que da la cota indicada. Los intervalos
encajados convergen a un punto y, por continuidad y cambio de signo, ese
punto es $M^*$. La cota del punto medio es geométrica.

Para la última afirmación, elíjase un entorno de $M^*$ sin otros quiebres.
Una vez que el encaje y el punto de evaluación están en ese entorno, cada
tramo lateral pertinente contiene la raíz en su clausura. La propuesta de
Newton es entonces exactamente $M^*$ por el
\cref{thm:terminacion-finita}(i), y pertenece al interior del encaje mientras
no se haya alcanzado la raíz. Si el contador obliga primero a una bisección,
la siguiente propuesta válida de Newton obtiene el mismo resultado. Esto
ocurre después de un número finito de pasos por la contracción del encaje.
\end{proof}

La bisección periódica permite recomendar este esquema por una garantía de
progreso explícita. Su costo efectivo se contrasta con el de las otras
alternativas en el \cref{cap:resultados}; la prueba no establece que sea el
único método del catálogo con garantía global.

\subsection{Los demás métodos abiertos}

\paragraph{Steffensen.}
$x_{k+1}=x_k-f(x_k)^2/\bigl(f(x_k+f(x_k))-f(x_k)\bigr)$. Alcanza orden $2$ sin
derivadas usando el propio valor de $f$ como incremento.

\begin{criticism}[Un fracaso instructivo]
\label{crit:steffensen}
Steffensen falla en el $100\%$ de los escenarios de este trabajo. La causa es
puramente dimensional: $f$ se mide en pesos de capital, con magnitud
$\sim10^8$, mientras que $x$ se mide en pesos de mesada, con magnitud
$\sim10^6$. El incremento auxiliar $x_k+f(x_k)$ resulta ser unas cien veces mayor
que la propia incógnita, sale del dominio admisible, es recortado al extremo
superior, y la diferencia del denominador se anula. El método no es «malo»: está
mal escalado. La variante \emph{escalada}, que usa $x_k+f(x_k)/L$ con
$L\approx\VP'$, converge en $\datoEvSteE$ evaluaciones. La lección es que el orden
de convergencia es invariante bajo cambio de unidades, pero la \emph{convergencia}
no lo es.
\end{criticism}

\paragraph{Müller.}
Interpola una parábola por los tres últimos puntos; orden $1{,}839$, la raíz real
de $t^3=t^2+t+1$ \citep{muller1956}.\footnote{El valor es $p\approx1{,}83929$, la llamada constante \emph{tribonacci}, que generaliza la razón áurea de la secante del mismo modo que Müller generaliza la interpolación lineal a la cuadrática. La sucesión de órdenes de los métodos de interpolación de grado creciente tiende a $2$ por abajo.} Sobre este problema pierde su ventaja: como
$f$ es afín en cada tramo, el coeficiente cuadrático de la parábola es
\emph{nulo} dentro de un mismo tramo y Müller degenera exactamente en la secante;
el término cuadrático solo se activa cuando los tres puntos abarcan un quiebre, es
decir, precisamente donde una parábola no describe bien a la función.

\paragraph{Halley.}
$x_{k+1}=x_k-2ff'/(2f'^2-ff'')$, de orden cúbico. Es el contraejemplo más nítido
del capítulo:

\begin{criticism}[El orden nominal alto no garantiza menor costo]
\label{crit:halley}
Halley es el método de mayor orden teórico del catálogo ($p=3$) y uno de los peores
en la práctica ($\datoEvHal$ evaluaciones, frente a $\datoEvNAN$ de Newton
analítico). La razón es estructural: el orden $3$ presupone $f\in C^3$ y explota la
información de $f''$\footnote{Halley publicó el método en 1694 como una mejora del de Newton, y su deducción original --- por aproximación mediante una hipérbola en lugar de una recta --- presupone explícitamente curvatura no nula. Sobre una función afín a trozos la hipérbola degenera en la recta, y con ella toda la ventaja del método.}. Pero aquí $f''\equiv0$ en casi todo punto, de modo que la
fórmula de Halley colapsa idénticamente en la de Newton --- sin ganancia alguna,
pero pagando tres evaluaciones por paso en lugar de una --- y en los quiebres la
segunda derivada numérica se comporta como $O(1/h)$ y desestabiliza el paso. Es la
ilustración más clara de la tesis metodológica de este trabajo: \emph{la elección
del método debe guiarse por la estructura del problema, no por el orden nominal}.
\end{criticism}

\paragraph{Punto fijo relajado.}
$x_{k+1}=x_k+\lambda f(x_k)$, con $\lambda<0$. Por el teorema de la aplicación
contractiva converge si $|1+\lambda f'|<1$, y el valor óptimo $\lambda=-1/f'$
reproduce exactamente Newton. En este problema $f'\approx\datoPendienteAfin$ es
prácticamente constante en todo el rango comercial, de modo que la elección
$\lambda=-1/\datoPendienteAfin$ es casi óptima \emph{globalmente} y el método de
orden $1$ rinde como uno de orden $2$: $\datoEvPF$ evaluaciones. Es un caso
notable en que la estructura afín permite «congelar» la derivada sin penalización.

\paragraph{Aitken--Steffensen.}
Acelera el punto fijo con el proceso $\Delta^2$ de Aitken. Convierte la
convergencia lineal en cuadrática al precio de dos evaluaciones por paso.

\section{Síntesis del capítulo}

El \cref{tab:prediccion} resume qué predice la teoría genérica y qué predice el
análisis estructural del \cref{cap:estructura} para cada método. El
\cref{cap:resultados} contrasta ambas columnas con el experimento.

\begin{table}[H]
\centering
\caption[Predicciones teóricas]{Lo que predice la teoría genérica frente a lo que
predice la estructura afín a trozos.}
\label{tab:prediccion}
\small
\begin{tabular}{@{}p{3.4cm}p{4.6cm}p{6.2cm}@{}}
\toprule
\textbf{Método} & \textbf{Teoría genérica} & \textbf{Predicción estructural} \\
\midrule
Bisección & lineal, razón $1/2$ & sin cambio: la cota no usa suavidad \\
Regula falsi & lineal & puede conservar un extremo; exacta si la cuerda queda en el tramo de la raíz \\
Secante & orden $\varphi$ & exacta con dos puntos en el tramo de la raíz (\cref{prop:secante-afin}) \\
Newton dif.\ fin. & orden $2$ & pendiente exacta si la ventana no cruza un quiebre; costo $3\times$ \\
Newton analítico & orden $2$ & un paso dentro del tramo de la raíz (\cref{cor:un-paso}) \\
Steffensen & orden $2$ & falla por escala (\cref{crit:steffensen}) \\
Müller & orden $1{,}839$ & degenera en secante \\
Halley & orden $3$ & degenera en Newton con triple costo \\
Punto fijo & lineal & casi óptimo por $f'$ constante \\
Newton semisuave & superlineal & terminación local exacta; alcance según el \cref{thm:terminacion-finita} \\
\bottomrule
\end{tabular}
\end{table}
